% ═══════════════════════════════════════════════════════════════
% MUSTERLÖSUNG: Atommodelle (Physik Klasse 10)
% Identisches Layout wie AB, Lösungen in ROT
% ═══════════════════════════════════════════════════════════════

\documentclass[10pt, a4paper]{article}

% ═══════════════════════════════════════════════════════════════
% PAKETE
% ═══════════════════════════════════════════════════════════════

\usepackage[utf8]{inputenc}
\usepackage[T1]{fontenc}
\usepackage[ngerman]{babel}
\usepackage{helvet}
\renewcommand{\familydefault}{\sfdefault}

\usepackage[margin=1.5cm, top=1.8cm, bottom=1.5cm]{geometry}
\usepackage{enumitem}
\usepackage{tabularx}
\usepackage{booktabs}
\usepackage{array}
\usepackage{multicol}

\usepackage{xcolor}
\usepackage{tcolorbox}
\usepackage{amssymb}
\usepackage{fancyhdr}
\usepackage{lastpage}
\usepackage[normalem]{ulem}

% Kompaktere Listen
\setlist{nosep, leftmargin=1.5em}

% ═══════════════════════════════════════════════════════════════
% FARBEN
% ═══════════════════════════════════════════════════════════════

\definecolor{physik}{HTML}{2c5aa0}
\definecolor{grau}{HTML}{888888}
\definecolor{hellgrau}{HTML}{f0f0f0}
\definecolor{success}{HTML}{2a7a4b}
\definecolor{fehler}{HTML}{c62828}
\definecolor{fehlerbg}{HTML}{ffebee}
\definecolor{loesung}{HTML}{c62828}

\newcommand{\fachfarbe}{physik}

% ═══════════════════════════════════════════════════════════════
% VARIABLEN
% ═══════════════════════════════════════════════════════════════

\newcommand{\thema}{Atommodelle}
\newcommand{\fach}{Physik}
\newcommand{\klasse}{Klasse 10}
\newcommand{\stunde}{Stunde 1-2}

% ═══════════════════════════════════════════════════════════════
% KOPF- UND FUSSZEILE
% ═══════════════════════════════════════════════════════════════

\pagestyle{fancy}
\fancyhf{}
\renewcommand{\headrulewidth}{0.4pt}
\renewcommand{\footrulewidth}{0pt}

\lhead{\small\fach{} \klasse{} -- \thema{} \textcolor{loesung}{\textbf{[MUSTERLÖSUNG]}}}
\rhead{\small Name: \rule{4cm}{0.4pt}}

\cfoot{\small\thepage/\pageref{LastPage}}

% ═══════════════════════════════════════════════════════════════
% BEFEHLE
% ═══════════════════════════════════════════════════════════════

% Lösung (rot)
\newcommand{\lsg}[1]{\textcolor{loesung}{\textbf{#1}}}

% Physik-Highlight
\newcommand{\ph}[1]{\textcolor{\fachfarbe}{\textbf{#1}}}

% Merksatz-Box
\newtcolorbox{merksatz}[1][]{
    colback=hellgrau,
    colframe=\fachfarbe,
    boxrule=0.8pt,
    arc=0pt,
    left=6pt, right=6pt, top=6pt, bottom=6pt,
    fonttitle=\bfseries\small,
    title=#1
}

% Fehler-Box
\newtcolorbox{fehlerbox}{
    colback=fehlerbg,
    colframe=fehler,
    boxrule=0.8pt,
    arc=0pt,
    left=6pt, right=6pt, top=4pt, bottom=4pt,
    fonttitle=\bfseries\small,
    title=Typische Fehler
}

% Abschnittstitel
\newcommand{\abschnitt}[1]{%
    \vspace{8pt}
    \noindent\textbf{\textcolor{\fachfarbe}{#1}}
    \vspace{4pt}
}

% ═══════════════════════════════════════════════════════════════
% DOKUMENT
% ═══════════════════════════════════════════════════════════════

\begin{document}

% ═══════════════════════════════════════════════════════════════
% TITEL
% ═══════════════════════════════════════════════════════════════

\begin{center}
{\large\bfseries Arbeitsblatt: \thema{} -- \stunde{} \textcolor{loesung}{[MUSTERLÖSUNG]}}\\[2pt]
{\small\textcolor{grau}{Begleitend zum Lernpfad -- in den Hefter übernehmen}}
\end{center}

\vspace{-2pt}
\hrule
\vspace{8pt}

% ═══════════════════════════════════════════════════════════════
% ZWEISPALTIG
% ═══════════════════════════════════════════════════════════════

\begin{multicols}{2}
\small

% ─────────────────────────────────────────────────────────────
% KASTEN 1: Die vier Atommodelle
% ─────────────────────────────────────────────────────────────

\abschnitt{Kasten 1: Die vier Atommodelle}

\begin{merksatz}[Historische Entwicklung]
Ergänze die Tabelle aus dem Lernpfad:

\vspace{6pt}

\renewcommand{\arraystretch}{1.6}
\begin{tabular}{|p{1.8cm}|p{0.8cm}|p{3.5cm}|}
\hline
\textbf{Modell} & \textbf{Jahr} & \textbf{Kernaussage} \\
\hline
Dalton & \lsg{1803} & \lsg{Unteilbare, massive Kugeln} \\
\hline
Thomson & \lsg{1897} & \lsg{Rosinenkuchen: + Masse mit e$^-$} \\
\hline
Rutherford & \lsg{1911} & \lsg{Kern-Hülle-Modell} \\
\hline
Bohr & \lsg{1913} & \lsg{Feste Bahnen (Energieniveaus)} \\
\hline
\end{tabular}
\renewcommand{\arraystretch}{1}
\end{merksatz}

\vspace{6pt}

% ─────────────────────────────────────────────────────────────
% KASTEN 2: Rutherford-Streuversuch
% ─────────────────────────────────────────────────────────────

\abschnitt{Kasten 2: Rutherford-Streuversuch}

\begin{merksatz}[Versuch und Ergebnisse]
\textbf{Aufbau:}\\
\lsg{Alpha}-Teilchen werden auf eine dünne\\
\lsg{Gold}-Folie geschossen.

\vspace{6pt}

\textbf{Beobachtungen $\rightarrow$ Schlussfolgerungen:}

\vspace{4pt}

\renewcommand{\arraystretch}{1.4}
\begin{tabular}{|p{2.8cm}|p{3.2cm}|}
\hline
\textbf{Beobachtung} & \textbf{Schlussfolgerung} \\
\hline
Die meisten gehen durch & \lsg{Atom ist größtenteils leer} \\
\hline
Wenige werden stark abgelenkt & \lsg{+ Ladung in kleinem Kern konzentriert} \\
\hline
\end{tabular}
\renewcommand{\arraystretch}{1}

\vspace{6pt}

\textbf{Größenverhältnis:}\\
Kerndurchmesser: $\approx$ \lsg{10$^{-15}$} m\\
Atomdurchmesser: $\approx$ \lsg{10$^{-10}$} m
\end{merksatz}

\vspace{6pt}

% ─────────────────────────────────────────────────────────────
% KASTEN 3: Kernaufbau
% ─────────────────────────────────────────────────────────────

\abschnitt{Kasten 3: Kernaufbau}

\begin{merksatz}[Bestandteile des Atoms]
\renewcommand{\arraystretch}{1.4}
\begin{tabular}{|l|c|c|c|}
\hline
\textbf{Teilchen} & \textbf{Ladung} & \textbf{Masse} & \textbf{Ort} \\
\hline
Proton & \lsg{+1} & \lsg{$\approx$1 u} & \lsg{Kern} \\
\hline
Neutron & \lsg{0} & \lsg{$\approx$1 u} & \lsg{Kern} \\
\hline
Elektron & \lsg{-1} & \lsg{$\approx$0} & \lsg{Hülle} \\
\hline
\end{tabular}
\renewcommand{\arraystretch}{1}

\vspace{8pt}

\textbf{Schreibweise:} $^{A}_{Z}$X

\vspace{4pt}

A = \lsg{Massenzahl} (Protonen + Neutronen)\\
Z = \lsg{Ordnungszahl} (Anzahl Protonen)\\
N = \lsg{Neutronenzahl} (Formel: \lsg{A - Z})
\end{merksatz}

% ═══════════════════════════════════════════════════════════
% SPALTE 2
% ═══════════════════════════════════════════════════════════

\columnbreak

% ─────────────────────────────────────────────────────────────
% ÜBUNG 1: Zuordnung (AFB I)
% ─────────────────────────────────────────────────────────────

\abschnitt{Übung 1: Zeitstrahl ergänzen}

Ordne die Jahreszahlen und Modelle chronologisch:

\vspace{4pt}

\textit{Jahreszahlen: 1803, 1897, 1911, 1913}

\vspace{6pt}

\begin{center}
\lsg{1803} $\rightarrow$ \lsg{1897} $\rightarrow$ \lsg{1911} $\rightarrow$ \lsg{1913}

\vspace{4pt}

\lsg{Dalton} $\rightarrow$ \lsg{Thomson} $\rightarrow$ \lsg{Rutherford} $\rightarrow$ \lsg{Bohr}
\end{center}

\vspace{6pt}

% ─────────────────────────────────────────────────────────────
% ÜBUNG 2: Lückentext (AFB II)
% ─────────────────────────────────────────────────────────────

\abschnitt{Übung 2: Rutherford erklären}

Ergänze den Lückentext:

\vspace{4pt}

Beim Streuversuch wurden \lsg{Alpha}-Teilchen

auf eine dünne Goldfolie geschossen. Die meisten

gingen \lsg{unabgelenkt} durch. Das zeigt:

Das Atom ist größtenteils \lsg{leer}.

Wenige Teilchen wurden stark \lsg{abgelenkt}.

Das zeigt: Die positive Ladung ist in einem

\lsg{kleinen} Kern konzentriert.

\vspace{8pt}

% ─────────────────────────────────────────────────────────────
% ÜBUNG 3: Kernaufbau (AFB II)
% ─────────────────────────────────────────────────────────────

\abschnitt{Übung 3: Teilchenzahlen bestimmen}

Bestimme für die folgenden Nuklide:

\vspace{6pt}

\begin{enumerate}[leftmargin=2em]
\item $^{12}_{6}$C: \quad p = \lsg{6} \quad n = \lsg{6} \quad e$^-$ = \lsg{6}

\vspace{2pt}

\item $^{235}_{92}$U: \quad p = \lsg{92} \quad n = \lsg{143} \quad e$^-$ = \lsg{92}

\vspace{2pt}

\item $^{56}_{26}$Fe: \quad p = \lsg{26} \quad n = \lsg{30} \quad e$^-$ = \lsg{26}
\end{enumerate}

\vspace{8pt}

% ─────────────────────────────────────────────────────────────
% ÜBUNG 4: Fehler begründen (AFB III)
% ─────────────────────────────────────────────────────────────

\abschnitt{Übung 4: Fehler erklären (AFB III)}

\begin{fehlerbox}
Erkläre, warum diese Aussagen falsch sind:

\vspace{6pt}

\begin{enumerate}[leftmargin=1.5em]
\item \textit{„Das Thomson-Modell wurde durch die Entdeckung des Elektrons widerlegt."}

\vspace{2pt}
\lsg{Falsch! Thomson entwickelte sein Modell WEGEN der Elektronenentdeckung. Es wurde durch den Rutherford-Versuch widerlegt.}

\vspace{6pt}

\item \textit{„Im Rutherford-Versuch wurden die meisten Teilchen stark abgelenkt."}

\vspace{2pt}
\lsg{Falsch! Die meisten gingen unabgelenkt durch. Nur sehr wenige (1:8000) wurden stark abgelenkt.}

\vspace{6pt}

\item \textit{„Die Massenzahl A gibt die Anzahl der Elektronen an."}

\vspace{2pt}
\lsg{Falsch! A = Protonen + Neutronen. Die Elektronenzahl entspricht bei neutralen Atomen der Protonenzahl Z.}
\end{enumerate}
\end{fehlerbox}

\vspace{8pt}

% ─────────────────────────────────────────────────────────────
% FREIES SCHREIBEN (AFB III)
% ─────────────────────────────────────────────────────────────

\abschnitt{Übung 5: Transfer}

Erkläre in 2-3 Sätzen: Warum musste nach dem Rutherford-Versuch ein neues Atommodell entwickelt werden?

\vspace{4pt}

\lsg{Das Kern-Hülle-Modell konnte nicht erklären, warum Elektronen nicht in den Kern stürzen (Energieverlust durch Strahlung). Außerdem konnte es die Linienspektren der Elemente nicht erklären. Bohr löste dies durch die Annahme fester Bahnen mit bestimmten Energieniveaus.}

\end{multicols}

% ═══════════════════════════════════════════════════════════════
% SEITE 2: SIMULATIONSPROTOKOLL
% ═══════════════════════════════════════════════════════════════

\newpage

\begin{center}
{\large\bfseries Simulationsprotokoll: Rutherford-Streuversuch \textcolor{loesung}{[MUSTERLÖSUNG]}}\\[2pt]
{\small\textcolor{grau}{PhET-Simulation: phet.colorado.edu/de/simulations/rutherford-scattering}}
\end{center}

\vspace{-2pt}
\hrule
\vspace{12pt}

% ─────────────────────────────────────────────────────────────
% VERSUCHSAUFBAU
% ─────────────────────────────────────────────────────────────

\abschnitt{Teil A: Versuchsaufbau verstehen}

\begin{merksatz}[Aufbau skizzieren]
Zeichne den Versuchsaufbau in die Box. Beschrifte: Strahlungsquelle, Alpha-Teilchen, Goldfolie, Zinksulfidschirm.

\vspace{6pt}
\lsg{Skizze: Links Strahlungsquelle (Radium) $\rightarrow$ Alpha-Teilchen (Pfeile) $\rightarrow$ Mitte: dünne Goldfolie (senkrecht) $\rightarrow$ Rundherum: Zinksulfidschirm (Halbkreis) mit Blitzen wo Teilchen auftreffen}
\vspace{40pt}

\end{merksatz}

\vspace{8pt}

% ─────────────────────────────────────────────────────────────
% VORHERSAGE (POE: Predict)
% ─────────────────────────────────────────────────────────────

\abschnitt{Teil B: Vorhersage (VOR der Simulation)}

\textbf{Wenn das Thomson-Modell stimmt} (gleichmäßig verteilte positive Ladung):

Was erwartest du? Kreuze an:

\vspace{4pt}

$\square$ Alle Teilchen werden stark abgelenkt

$\square$ Alle Teilchen gehen unabgelenkt durch

\lsg{$\boxtimes$ Alle Teilchen werden gleichmäßig leicht abgelenkt}

\vspace{6pt}

Begründe deine Vorhersage: \lsg{Bei gleichmäßig verteilter Ladung wirken nur schwache Kräfte auf die Alpha-Teilchen, daher nur leichte Ablenkung.}

\vspace{8pt}

% ─────────────────────────────────────────────────────────────
% BEOBACHTUNG (POE: Observe)
% ─────────────────────────────────────────────────────────────

\abschnitt{Teil C: Beobachtung (WÄHREND der Simulation)}

Starte die Simulation und beobachte. Zähle bei 100 abgefeuerten Teilchen:

\vspace{6pt}

\begin{tabular}{|p{6cm}|c|c|}
\hline
\textbf{Beobachtung} & \textbf{Anzahl} & \textbf{Anteil} \\
\hline
Teilchen gehen unabgelenkt durch & \lsg{$\approx$90} & \lsg{$\approx$90} \% \\
\hline
Teilchen werden leicht abgelenkt (<45°) & \lsg{$\approx$9} & \lsg{$\approx$9} \% \\
\hline
Teilchen werden stark abgelenkt (>45°) & \lsg{$\approx$1} & \lsg{$\approx$1} \% \\
\hline
Teilchen werden zurückgeworfen (>90°) & \lsg{$<$1} & \lsg{$<$1} \% \\
\hline
\end{tabular}

\vspace{8pt}

Stimmt deine Vorhersage mit der Beobachtung überein? $\square$ Ja \quad \lsg{$\boxtimes$ Nein}

\vspace{8pt}

% ─────────────────────────────────────────────────────────────
% ERKLÄRUNG (POE: Explain)
% ─────────────────────────────────────────────────────────────

\abschnitt{Teil D: Erklärung (NACH der Simulation)}

\textbf{1. Warum gehen die meisten Teilchen unabgelenkt durch?}

\vspace{4pt}
\lsg{Das Atom ist größtenteils leerer Raum. Die Alpha-Teilchen treffen auf nichts und fliegen ungehindert durch.}

\vspace{8pt}

\textbf{2. Warum werden manche Teilchen stark abgelenkt oder zurückgeworfen?}

\vspace{4pt}
\lsg{Diese Teilchen kommen dem kleinen, positiven Kern sehr nahe. Die starke elektrostatische Abstoßung (+ gegen +) lenkt sie stark ab.}

\vspace{8pt}

\textbf{3. Was bedeutet das für den Aufbau des Atoms?}

\vspace{4pt}
\lsg{Das Atom besteht aus einem winzigen, positiven Kern (Durchmesser $\approx 10^{-15}$ m) und einer riesigen, fast leeren Hülle (Durchmesser $\approx 10^{-10}$ m). Die gesamte positive Ladung und fast die gesamte Masse ist im Kern konzentriert.}

\vspace{12pt}

% ─────────────────────────────────────────────────────────────
% ZUSATZEXPERIMENT
% ─────────────────────────────────────────────────────────────

\abschnitt{Teil E: Zusatzexperiment (für Schnelle)}

Verändere in der Simulation die \textbf{Anzahl der Protonen} im Kern.

\vspace{6pt}

Was passiert bei \textbf{mehr} Protonen? \lsg{Mehr Teilchen werden abgelenkt, stärkere Ablenkwinkel.}

\vspace{6pt}

Erkläre, warum: \lsg{Mehr Protonen = größere positive Ladung = stärkere elektrostatische Abstoßung der Alpha-Teilchen.}

\end{document}
